% Options for packages loaded elsewhere
\PassOptionsToPackage{unicode}{hyperref}
\PassOptionsToPackage{hyphens}{url}
\PassOptionsToPackage{dvipsnames,svgnames,x11names}{xcolor}
%
\documentclass[
  10pt,
]{article}
\usepackage{amsmath,amssymb}
\usepackage{iftex}
\ifPDFTeX
  \usepackage[T1]{fontenc}
  \usepackage[utf8]{inputenc}
  \usepackage{textcomp} % provide euro and other symbols
\else % if luatex or xetex
  \usepackage{unicode-math} % this also loads fontspec
  \defaultfontfeatures{Scale=MatchLowercase}
  \defaultfontfeatures[\rmfamily]{Ligatures=TeX,Scale=1}
\fi
\usepackage{lmodern}
\ifPDFTeX\else
  % xetex/luatex font selection
  \setmainfont[]{DejaVu Serif}
  \setsansfont[]{DejaVu Sans}
  \setmonofont[]{DejaVu Sans Mono}
  \setmathfont[]{Latin Modern Math}
\fi
% Use upquote if available, for straight quotes in verbatim environments
\IfFileExists{upquote.sty}{\usepackage{upquote}}{}
\IfFileExists{microtype.sty}{% use microtype if available
  \usepackage[]{microtype}
  \UseMicrotypeSet[protrusion]{basicmath} % disable protrusion for tt fonts
}{}
\makeatletter
\@ifundefined{KOMAClassName}{% if non-KOMA class
  \IfFileExists{parskip.sty}{%
    \usepackage{parskip}
  }{% else
    \setlength{\parindent}{0pt}
    \setlength{\parskip}{6pt plus 2pt minus 1pt}}
}{% if KOMA class
  \KOMAoptions{parskip=half}}
\makeatother
\usepackage{xcolor}
\usepackage[margin=2.0cm]{geometry}
\usepackage{longtable,booktabs,array}
\usepackage{calc} % for calculating minipage widths
% Correct order of tables after \paragraph or \subparagraph
\usepackage{etoolbox}
\makeatletter
\patchcmd\longtable{\par}{\if@noskipsec\mbox{}\fi\par}{}{}
\makeatother
% Allow footnotes in longtable head/foot
\IfFileExists{footnotehyper.sty}{\usepackage{footnotehyper}}{\usepackage{footnote}}
\makesavenoteenv{longtable}
\setlength{\emergencystretch}{3em} % prevent overfull lines
\providecommand{\tightlist}{%
  \setlength{\itemsep}{0pt}\setlength{\parskip}{0pt}}
\setcounter{secnumdepth}{-\maxdimen} % remove section numbering
% doc_header.tex -- LaTeX preamble additions used by build_docs.sh for every
% Markdown -> LaTeX -> PDF conversion (supplements, README, submission files).
\usepackage{etoolbox}
% longtable may not be loaded by pandoc when a document has no table
\usepackage{longtable}
% bold vectors: \bm -> \symbf under unicode-math (XeLaTeX), bm package otherwise
\makeatletter
\@ifpackageloaded{unicode-math}{\providecommand{\bm}[1]{\symbf{#1}}}{\usepackage{bm}}
\makeatother
% smaller type in long tables so that wide audit tables fit the text width;
% inside a table a line break is allowed after an underscore of plain text
% (pandoc writes it as \_), so that long identifiers wrap inside their column
% rows are set a little apart, so that symbols with indices do not touch
\AtBeginEnvironment{longtable}{\footnotesize\renewcommand{\arraystretch}{1.45}%
  \let\sptunderscore\_\def\_{\sptunderscore\allowbreak}}
\setlength{\LTcapwidth}{\textwidth}
% allow line breaks after slashes and hyphens inside table cells
\usepackage{xurl}
% code blocks: break over-long lines instead of running into the margin.
% Plain ``` blocks are typeset by pandoc as {verbatim}, highlighted ones
% (```bash) as {Highlighting}; both are re-declared as breakable fancyvrb
% environments.  The continuation line is marked with a small hooked arrow.
\usepackage{fvextra}
\RecustomVerbatimEnvironment{verbatim}{Verbatim}{breaklines,breakanywhere,fontsize=\small}
\makeatletter
\@ifundefined{Highlighting}{}{%
  \RecustomVerbatimEnvironment{Highlighting}{Verbatim}{breaklines,breakanywhere,fontsize=\small,commandchars=\\\{\}}}
\makeatother
% inline code (pandoc writes it as \texttt{...}): never hyphenate -- an inserted
% hyphen could be read as part of a file name -- but allow a line break after
% an underscore, slash, comma, colon or equals sign, so that long paths and
% column lists wrap inside the text width and inside table columns.
% Inside a formula (e.g. the name of a notebook symbol in an equation) the same
% command must first leave math mode, which \text of amsmath does.
\ExplSyntaxOn
\tl_new:N \l__spt_code_tl
\cs_new_protected:Npn \__spt_texttt:n #1
  {
    \tl_set:Nn \l__spt_code_tl {#1}
    \regex_replace_all:nnN { ([\/\,\:\=]) } { \1 \c{allowbreak} } \l__spt_code_tl
    \regex_replace_all:nnN { \c{_} } { \c{_} \c{allowbreak} } \l__spt_code_tl
    \group_begin: \ttfamily \hyphenchar \font = -1 \scan_stop: \l__spt_code_tl \group_end:
  }
\RenewDocumentCommand \texttt { m }
  {
    \mode_if_math:TF { \text { \__spt_texttt:n {#1} } } { \__spt_texttt:n {#1} }
  }
\ExplSyntaxOff
% A figure is written in the Markdown source as an image paragraph followed by
% a caption paragraph ("**Figure S1.** ..."), so that the caption is also
% visible on GitHub.  An image that starts a paragraph is centered.
\usepackage{graphicx}% pandoc loads it only for documents that contain an image
\NewCommandCopy\sptincludegraphics\includegraphics
\RenewDocumentCommand\includegraphics{O{} m}{%
  \ifvmode\begin{center}\sptincludegraphics[#1]{#2}\end{center}%
  \else\sptincludegraphics[#1]{#2}\fi}
% keep figures where they are referenced
\usepackage{float}
\floatplacement{figure}{H}
% running header with the document title
\usepackage{fancyhdr}
\pagestyle{fancy}
\fancyhf{}
\fancyhead[L]{\small PROJECT spacetimes, Revision\_Deflation --- documentation and provenance}
\fancyhead[R]{\small\thepage}
\renewcommand{\headrulewidth}{0.4pt}
% 16x16 matrices: tighter columns so that a full matrix fits the text width
\setlength{\arraycolsep}{3.2pt}
% The documents of this project contain long tagged symbols, file paths and
% identifiers that cannot be hyphenated.  A paragraph is therefore allowed to
% have looser lines before TeX gives up and lets a line stick out of the margin.
\AtBeginDocument{\tolerance=2000 \emergencystretch=5em \hbadness=2000}
\ifLuaTeX
  \usepackage{selnolig}  % disable illegal ligatures
\fi
\IfFileExists{bookmark.sty}{\usepackage{bookmark}}{\usepackage{hyperref}}
\IfFileExists{xurl.sty}{\usepackage{xurl}}{} % add URL line breaks if available
\urlstyle{same}
\hypersetup{
  pdftitle={RP9 --- Error log of Revision\_Deflation},
  pdfauthor={Prepared for Patrick Lee Nash, Ph.D.},
  colorlinks=true,
  linkcolor={blue},
  filecolor={Maroon},
  citecolor={blue},
  urlcolor={blue},
  pdfcreator={LaTeX via pandoc}}

\title{RP9 --- Error log of Revision\_Deflation}
\usepackage{etoolbox}
\makeatletter
\providecommand{\subtitle}[1]{% add subtitle to \maketitle
  \apptocmd{\@title}{\par {\large #1 \par}}{}{}
}
\makeatother
\subtitle{Provenance of PROJECT spacetimes, Revision\_Deflation: every
error made during the work on the revision, how it was found, and how it
was corrected}
\author{Prepared for Patrick Lee Nash, Ph.D.}
\date{November 13, 2026}

\begin{document}
\maketitle

\hypertarget{purpose}{%
\section{Purpose}\label{purpose}}

Hard Rule 7 of the task specification forbids shortcuts and demands that
problems be explained. As the provenance P9 does for the main project,
this file records every error that was made while the revision was
carried out --- including errors that were corrected before anything was
committed. For each it states what was wrong, how it was detected, what
was done, and where the corrected statement can be checked. An error
that is not listed here and is found later is to be added, not silently
repaired.

The errors are numbered R-E1, R-E2, \ldots{} in the order in which they
were recorded.

\hypertarget{errors-of-physics-and-of-mathematics}{%
\section{1. Errors of physics and of
mathematics}\label{errors-of-physics-and-of-mathematics}}

\begin{longtable}[]{@{}
  >{\raggedright\arraybackslash}p{(\columnwidth - 8\tabcolsep) * \real{0.0435}}
  >{\raggedright\arraybackslash}p{(\columnwidth - 8\tabcolsep) * \real{0.3043}}
  >{\raggedright\arraybackslash}p{(\columnwidth - 8\tabcolsep) * \real{0.2319}}
  >{\raggedright\arraybackslash}p{(\columnwidth - 8\tabcolsep) * \real{0.2464}}
  >{\raggedright\arraybackslash}p{(\columnwidth - 8\tabcolsep) * \real{0.1739}}@{}}
\toprule\noalign{}
\begin{minipage}[b]{\linewidth}\raggedright
No.
\end{minipage} & \begin{minipage}[b]{\linewidth}\raggedright
What was wrong
\end{minipage} & \begin{minipage}[b]{\linewidth}\raggedright
How it was found
\end{minipage} & \begin{minipage}[b]{\linewidth}\raggedright
Correction
\end{minipage} & \begin{minipage}[b]{\linewidth}\raggedright
Where to check
\end{minipage} \\
\midrule\noalign{}
\endhead
\bottomrule\noalign{}
\endlastfoot
R-E1 & For unequal rates (7+1) it was first asserted that the coordinate
that becomes energetic first decides the number of created particles
whenever the two crossings are at least three e-foldings apart. This is
false when the second direction deflates more slowly. & The check failed
for a slower second direction. & The criterion is that the second
physical wavenumber is still below \(e^{-3}E_\perp\) at the first
crossing. & check \texttt{71d-M3} \\
R-E2 & It was first asserted that all eight bounded solutions of the 4+4
field fall by the same factor. This holds only when one rate governs. &
The check failed for three unequal rates. & With unequal rates the
factor depends on the internal state (spread 0.40 at
\(S_{\rm WKB}=30\)); stated so. & check \texttt{44d-M3} \\
R-E5 & The critical value of \(C\) at which matter becomes abnormal was
required to agree with the estimate \(\pi^2m^2/4k_F^2\) of the main
project within 60 per cent at all densities. & The check failed at low
density. & The estimate is reached only at high density; at
\(k_F=0.2\,m\) the critical value is 2.46 times larger. & check
\texttt{def-H5}; RS4, Sec. 3.1 \\
R-E6 & The effect of the loss of mass of bound pairs on the Hubble
constant was first asserted as ``at most 2 km/s/Mpc'' --- a guess ---
and computed with \(\alpha_{VS}=3\), where the formula for the binding
energy is not valid. & Comparison of the assertion with the computed
numbers; review of the range of validity. & Examples with
\(\alpha_{VS}\le1\); at constant \(G_N\) the gain is at most \(+0.48\)
km/s/Mpc. & check \texttt{def-C4}; RS3, Sec. 4 \\
R-E7 & A bound on the rate was derived from the constancy of \(m_p/m_e\)
by assuming that the electron mass follows the four-dimensional vacuum
value of the Vorton-Scalar. Under that assumption the mass is the
product of a coupling that grows and a vacuum value that falls, and is
constant. & Re-derivation of the reduction of the couplings. & The row
was replaced by the bound from the Fermi constant at fixed masses. &
check \texttt{def-C5}; RS3, Sec. 3 \\
R-E9 & It was asserted that the amplitude \(\sigma_8\) of clustering
falls when Newton's constant grows with the deflation. & The computed
number has the other sign when the Hubble constant is the one that keeps
the sound-horizon angle. & Both cases are stated: \(-0.317\) at fixed
Hubble constant, \(+0.102\) at fixed angle. (These two numbers were
themselves superseded: see R-E33.) & check \texttt{def-D4}; RS3, Sec.
5 \\
R-E11 & The coherence length of the Ginzburg--Landau energy was computed
with a coefficient too large by the number of internal states, a factor
2 in the length. & Comparison with the weak-coupling value
\(0.1333\,v_F/T_c\). & Corrected; the weak-coupling formulas are stated
to hold only for \(T_c<0.05\,E_F\). & check \texttt{def-P3}; RS4, Sec.
2.3 \\
R-E12 & The law ``the condensation energy has twice the logarithmic
slope of the gap'' was asserted for all moments; and a counter of
strongly coupled moments mixed the two particles. & The check failed at
the strongly coupled moments. & The law is stated for weak coupling; the
counter is per particle. & check \texttt{def-P1} \\
R-E13 & The text of a check said that the fermion mass in self-bound
matter ``falls to'' zero along the clock. It is zero at every moment
after the threshold. & Reading the data file. & Text corrected. & check
\texttt{def-P7} \\
R-E14 & A routine of the main project that finds the melting temperature
treats ``the Hartree potential is not defined at 1.05 times the vacuum
value'' as ``melted''. For \(m_s/m=0.3\) and \(C\ge32\) this produced
temperatures that are artefacts; a check of monotonic decrease had
passed with them. & A jump in the figure. & The program marks such
moments as not defined (26 of 30 are defined). The table of the main
project, for \(C\le9\), is not affected. & check \texttt{def-H7}; RS4,
Sec. 4 \\
R-E17 & The table of exponents listed \(G_{VS}\) as \(g_{VS}^2/m_s^2\);
the main registry defines it as \(1/(\sqrt2\,v_{VS}^2)\) (the exponent
is the same). & Comparison with the main registry. & Corrected. &
\texttt{data/coupling\_exponents.csv} \\
R-E18 & The independent Python script had the wrong sign of one term in
the three special Killing vectors of the case of equal rates. The Rust
program was right. & The independent check failed. & Sign corrected,
with the derivation in a comment. & check \texttt{I44d-G4},
\texttt{I71d-G4} \\
R-E21 & The first table of the growth of the gap multiplied the
logarithmic slope by the pairing strength of the radial mode alone and
found 0.96 and 0.77 instead of 1. & Disagreement with check
\texttt{def-H3}, which uses the sum of the strengths of both modes. &
The table uses the pairing strength of the vector channel,
\(\lambda_S+\lambda_P\); the check says so explicitly. & RS1, Sec.
3.2 \\
R-E23 & \textbf{The label H2 was attached to the wrong hypothesis.} In
the main project H2 is the author's hypothesis that the superfluid state
is the Higgs field. In the programs and first documents of the revision
the label was used for a different assumption --- that the vacuum value
of the Vorton-Scalar sets the electroweak scale. The actual consequence
of H2 on a deflating spacetime had therefore not been examined. &
Reading the supplement S2 of the main project while writing RS2. & The
other assumption is now called HV. The consequence of H2 was derived and
computed: the electroweak scale follows \(\alpha_{VS}\), and atomic
clocks bound the rate to \(1.16\times10^{-6}H_0/n_{\rm EW}\). & check
\texttt{def-C6}; RS2, Sec. 5; RS3, Secs. 1, 3 \\
R-E25 & The statement of a check, and with it the supplement RS1 and the
table of RP0, spoke of ``90 pair channels''. The program had counted
each channel once for each of the three momenta: there are 15
representative channels on each spacetime, 30 in all, evaluated at three
momenta and 13 angles (1170 numbers). The comparison itself was complete
and correct. & Comparison with the data file of the main project while
writing the manuscripts. & The program counts the distinct channels; the
statement and the documents say 15 on each spacetime. & check
\texttt{def-K1}; RS1, Sec. 3.1 \\
R-E26 & The locally measured Hubble constant was written as ``73'' in
two statements of the program, in a figure and in RS3, and the distance
of the model from it as 5.6 km/s/Mpc. The main project compares with
\(73.50\pm0.81\) km/s/Mpc; the distance from 67.36 is 6.1. No conclusion
changes (the gain of at most \(+0.48\) remains below a tenth of it). &
Comparison with the manuscripts of the main project while writing the
manuscripts of the revision. & The program takes the value 73.50 and its
uncertainty as stated; statements, figure and RS3 corrected. & checks
\texttt{def-C3}, \texttt{def-C4}; \texttt{figures/fig\_hubble} \\
R-E28 & The summary of RS3 gave the temperature of the quanta that the
deflation creates as ``\(10^{-33}\) K at the allowed rate''. That is its
value at the bound of lunar laser ranging (\(2.9\times10^{-33}\) K); at
the largest rate that any of the three bounds allows it is
\(2.1\times10^{-31}\) K. The conclusion (nothing is created today) is
not affected. & The check of the numbers of the running text
(\texttt{check\_numbers.py}), written for the manuscripts, found that
the number was stated by no check. & The text gives both values; check
\texttt{def-D6} states them. & RS3, Sec. 0 and 7; check
\texttt{def-D6} \\
R-E30 & The caption of the table of states of the order parameter in the
drafts of the manuscripts gave the unit of the two gaps of a state as
``the square of the magnitude of the order parameter''. For the
half-paired state the tabulated value 1 would then have to be 2. The
numbers are in units of \(\Delta_F^2\), the square of the gap of a
unitary state, and each state is taken at the magnitude at which its
potential is stationary; the data file says so, and RS1 gave no unit.
The numbers themselves were right. & Re-reading the pages of the
finished manuscript against the identity for the gap matrix. & Caption
corrected; RS1 states the units. & manuscripts, Sec. X; RS1, Sec. 5.2;
\texttt{data/order\_parameter\_states.csv} \\
R-E33 & \textbf{The growth equation of RS3 was not the limit of the
eight-dimensional theory.} RS3, Sec. 5, the manuscripts and the check
\texttt{def-D4} integrated the familiar equation
\(d^2\delta_m/dx4^2+2H\,d\delta_m/dx4=4\pi G_N(x4)\rho_m\delta_m\) and
called it the Newtonian limit of the theory with assumptions A1 and A2.
It is not. With unperturbed extra scale factors the eight-dimensional
equations give a friction reduced by \(\mathrm{gsum}/2\) and, depending
on what is assumed about the agent that keeps the rates constant, a
reduced source as well. The quoted results (\(-0.317\) and \(+0.102\)
per unit of \(\mathrm{gsum}/H_0\) for the amplitude; factors 1.0027 and
1.0589 for the growth rate) belong to neither case. & The derivation of
the perturbation equations for RS5 by computer algebra; then the
complete Einstein--Boltzmann calculation, whose growth disagreed with
that equation by a quarter of \(\mathrm{gsum}/H_0\). & Two closures are
computed and a third analysed; RS3, Sec. 5, the manuscripts and the
check state the results of each and that the growth is not a prediction
of the hypotheses. The routine keeps the four-dimensional equation for
comparison only. & checks \texttt{def-D4}, \texttt{def-E8},
\texttt{Idef-P4} to \texttt{Idef-P7}; RS5, Sec. 4; RS3, Sec. 5 \\
R-E34 & \textbf{The fall of the Hubble constant was underestimated by
six per cent.} The command \texttt{cosmology} held the redshift of last
scattering and the helium fraction at their values without deflation
when it kept the angle of the sound horizon. With weaker gravity in the
past, last scattering is later and the helium fraction lower, and both
enlarge the sound horizon further. The slope was given as \(-0.911\)
instead of \(-0.969\), the values for \(\mathrm{gsum}=0.01\), 0.03,
\(0.1\,H_0\) as 66.74, 65.48, 60.77 instead of 66.70, 65.36, 60.37
km/s/Mpc, and the changes at the bounds of lunar laser ranging and
nucleosynthesis as \(-0.199\) and \(-4.13\) instead of \(-0.211\) and
\(-4.38\) km/s/Mpc. No conclusion changes; the fall is larger. & The
copy of the recombination program, run on the deflating background for
RS5, moved last scattering. & \texttt{cosmology} computes the full
thermal history and tabulates the old approximation beside it; an
independent recombination code confirms both (\texttt{Idef-C2}). The
examples with a varying rest mass (RS3, Sec. 4) still use the
approximation, and say so. & checks \texttt{def-C2}, \texttt{def-C3},
\texttt{def-B2}, \texttt{Idef-C1}, \texttt{Idef-C2}; RS3, Sec. 2; RS5,
Sec. 6 \\
R-E35 & The earlier version of RP0, Sec. 5, and of RS3, Sec. 9, named as
the simplest assumption for the perturbations that ``the scale factors
of \(x5,x6,x7\) are not perturbed, and the perturbations of ordinary
spacetime obey the four-dimensional equations with the Newton's constant
of the moment''. The second half does not follow from the first and is
not a consistent system: the eight-dimensional equations contain
\(H_e=H_{\rm cf}-a_{\rm sp}\mathrm{gsum}/2\) in place of the conformal
Hubble rate, and of the four scalar Einstein equations only two can hold
with matter alone. & The derivation by computer algebra; and in the
first copy of the Einstein--Boltzmann program the equation of the shear,
tested numerically along a solution, failed at the first order in the
rate. & Assumption A3 with its closures (RS5, Secs. 1 to 3); every
result is marked as depending on the closure or not. & checks
\texttt{I44d-P1}, \texttt{I71d-P1}, \texttt{Idef-P3}, \texttt{def-E4},
\texttt{def-E5}; RP0, Sec. 5 \\
R-E38 & A unit test of the growth equations was first written with
guessed first-order coefficients of the growth rate (\(-1/10\),
\(-1/2\), \(-3/5\) times \(\mathrm{gsum}/H\)). The test failed. The
coefficients, derived by a series solution, are \(-9/16\), \(-13/16\)
and \(-11/16\); and the test model needed a vanishing cosmological term,
which otherwise shifts the numbers by 0.006. & The unit test. & Derived,
confirmed by computer algebra, and tested to \(10^{-4}\). & unit test
\texttt{growth\_with\_deflation\_follows\_the\_first\_order\_laws};
check \texttt{Idef-P4} \\
R-E43 & The response of the helium fraction to the smaller Newton's
constant at nucleosynthesis was computed as \(\tfrac{43}{7}(e^{s}-1)\)
equivalent neutrino species for every model. That is right for the
standard radiation content only. In a model with extra radiation of the
density \(\Delta N\) during nucleosynthesis the expansion rate is
proportional to \(e^{s}(\tfrac{43}{7}+\Delta N)\), and the equivalent
number is \((\tfrac{43}{7}+\Delta N)(e^{s}-1)\). The difference is of
the second order (the product of the rate and of \(\Delta N\)):
\(4\times10^{-5}\) in the helium fraction for \(\Delta N=1\) at the rate
that lunar laser ranging allows, and \(4\times10^{-4}\) for
\(\Delta N=0.5\) at the rate that nucleosynthesis allows, where it
changed the fourth decimal of two entries of a table. & Re-reading the
routine while the first run of the command \texttt{families} was under
way. & The routine uses the exact correspondence; a unit test checks it;
the commands \texttt{background} and \texttt{families} were run again.
Models without extra radiation are not affected, to the last bit. &
\texttt{eb\_cosmo.rs}; unit test
\texttt{newtons\_constant\_at\_nucleosynthesis\_lowers\_the\_helium\_fraction};
RS5, Sec. 5.1 \\
\end{longtable}

\hypertarget{errors-of-numerical-method}{%
\section{2. Errors of numerical
method}\label{errors-of-numerical-method}}

\begin{longtable}[]{@{}
  >{\raggedright\arraybackslash}p{(\columnwidth - 8\tabcolsep) * \real{0.0435}}
  >{\raggedright\arraybackslash}p{(\columnwidth - 8\tabcolsep) * \real{0.3043}}
  >{\raggedright\arraybackslash}p{(\columnwidth - 8\tabcolsep) * \real{0.2319}}
  >{\raggedright\arraybackslash}p{(\columnwidth - 8\tabcolsep) * \real{0.2464}}
  >{\raggedright\arraybackslash}p{(\columnwidth - 8\tabcolsep) * \real{0.1739}}@{}}
\toprule\noalign{}
\begin{minipage}[b]{\linewidth}\raggedright
No.
\end{minipage} & \begin{minipage}[b]{\linewidth}\raggedright
What was wrong
\end{minipage} & \begin{minipage}[b]{\linewidth}\raggedright
How it was found
\end{minipage} & \begin{minipage}[b]{\linewidth}\raggedright
Correction
\end{minipage} & \begin{minipage}[b]{\linewidth}\raggedright
Where to check
\end{minipage} \\
\midrule\noalign{}
\endhead
\bottomrule\noalign{}
\endlastfoot
R-E3 & The test of the order of convergence of the integrator compared
the result with the exact value at a \emph{finite} end point, where the
two differ by the residual of the projection, not by the error of
integration. & The apparent order was wrong. & The test compares runs
with different steps. & unit test
\texttt{the\_integrator\_converges\_with\_order\_four} \\
R-E8 & The number of created scalar quanta was read off where the
physical wavenumber was 300 times the frequency; the residual,
\(6\times10^{-4}\), was large compared with the smallest numbers. & The
deviations of the check. & Read off at 3000 times: \(1.9\times10^{-5}\).
& check \texttt{def-B3} \\
R-E10 & The ``epoch of strong self-interaction'' was first defined by
the threshold 1 cm\(^2\)/g, which the example particle exceeds at all
times. & The interval came out as the whole range. & Threshold 100
cm\(^2\)/g. & check \texttt{def-D3} \\
R-E36 & The closure ``Shear'' was first built into the copy of the
Einstein--Boltzmann program, with the equation of the shear integrated
as an equation of motion. In that form the equations are stiff: the
momentum density of the agent relaxes at the rate \(k^2/3H_{\rm cf}\).
The explicit integrator of the main project could not follow it; even at
zero rate the temperature spectrum differed from that of the main
project by \(1.6\times10^{-3}\), and the sources at large wavenumbers
and late times were wrong by more than their size. & The test that the
copy reproduces the main project at zero rate. & The closure was removed
from the Rust program. It is treated analytically (RS5, Eq. (3.5):
inside the horizon it is the closure ``Pressure'') and integrated for
pressureless matter with an implicit method in the independent script. &
RS5, Sec. 3.2; check \texttt{Idef-P6} \\
R-E37 & A second design computed the derivatives of the metric
potentials that the source of the temperature needs by numerical
differences along the Runge--Kutta solution. The background of the
program is a table with linear interpolation; second derivatives through
it are meaningless, and the source had the wrong sign at one of the
tested times. & A test against the analytic expression at a negligible
rate. & The gaps are computed from the two identities of RS5, Sec. 3.1,
as algebraic expressions; only the derivative of the (small) shear gap
is taken by a difference, of first order and along the solution. &
\texttt{eb\_boltzmann.rs}; unit tests of \texttt{spectra\_defl.rs};
check \texttt{def-E5} \\
R-E39 & Two tests of the new program were themselves wrong. (i) The test
at a ``negligible'' rate used \(10^{-12}\), at which the adaptive
integrator chooses other steps and the spectra differ at \(10^{-9}\);
the rate \(10^{-30}\) takes the same branches and gives identical bits.
(ii) The gap of the equation of the shear was first compared with
\(a_{\rm sp}\mathrm{gsum}\) times a potential, which has the wrong
dimension; the measure is \((a_{\rm sp}\mathrm{gsum}/H_{\rm cf})\) times
the terms of the equation. & The tests failed, or gave numbers that
changed with the unit of length. & Corrected as stated. & unit tests of
\texttt{spectra\_defl.rs}; check \texttt{def-E5} \\
\end{longtable}

\hypertarget{errors-of-notation-hard-rule-5}{%
\section{3. Errors of notation (Hard Rule
5)}\label{errors-of-notation-hard-rule-5}}

\begin{longtable}[]{@{}
  >{\raggedright\arraybackslash}p{(\columnwidth - 8\tabcolsep) * \real{0.0435}}
  >{\raggedright\arraybackslash}p{(\columnwidth - 8\tabcolsep) * \real{0.3043}}
  >{\raggedright\arraybackslash}p{(\columnwidth - 8\tabcolsep) * \real{0.2319}}
  >{\raggedright\arraybackslash}p{(\columnwidth - 8\tabcolsep) * \real{0.2464}}
  >{\raggedright\arraybackslash}p{(\columnwidth - 8\tabcolsep) * \real{0.1739}}@{}}
\toprule\noalign{}
\begin{minipage}[b]{\linewidth}\raggedright
No.
\end{minipage} & \begin{minipage}[b]{\linewidth}\raggedright
What was wrong
\end{minipage} & \begin{minipage}[b]{\linewidth}\raggedright
How it was found
\end{minipage} & \begin{minipage}[b]{\linewidth}\raggedright
Correction
\end{minipage} & \begin{minipage}[b]{\linewidth}\raggedright
Where to check
\end{minipage} \\
\midrule\noalign{}
\endhead
\bottomrule\noalign{}
\endlastfoot
R-E4 & The author's names \texttt{g55}, \texttt{g66}, \texttt{g77} for
the rates were replaced by a Greek letter with an index, on the ground
that they resemble the components of the metric. & The author objected.
& The author's names are used exactly; the components of the metric are
distinguished by their subscripts and by the rule of RP1, Sec. 2. & RP1;
check RS4 \\
R-E16 & Names that the main registry uses for other objects were used
for new objects in data files and statements (among them \texttt{tau},
\texttt{omega\_c}, \texttt{mu}, \texttt{nu}, \texttt{kappa},
\texttt{eps}, \texttt{Q}, \texttt{M}, \texttt{G}, \texttt{D},
\texttt{f}, \texttt{u}, \texttt{A}, \texttt{P}). During the renaming, a
first rule would also have changed the identifier of a check; it was
restricted before it was applied. & Systematic comparison with the main
registry and with the symbols of the notebooks. & All renamed; six
machine checks now guard the registry. & RP1 \\
R-E19 & The first draft of DMD used unregistered letters
(\(\varepsilon\), \(\nu\), \(N\), \(C\), \(z\), \(Y\), \(\xi\), \(E_0\))
and the program name \texttt{h0}, which is a symbol of a notebook; and
it called the affine parameter the proper time ``multiplied by'' the
mass (it is divided by the mass). & Re-reading against the registries;
check RS3. & Replaced by registered symbols before the first commit of
the document. & DMD; RP1 \\
R-E20 & The renaming of R-E16 replaced the gap matrix \texttt{M} by
\texttt{m\_bos}, the name of a boson mass, in two statements. & Reading
the output while writing RS1. & Restored as \texttt{M\_gap}. & check
\texttt{def-K3} \\
R-E22 & The scalar mode and its exponents were called \texttt{S} and
\texttt{sigma} in headers of data files and in a figure; both letters
have other meanings in the main registry. & The same reading. &
\texttt{phi\_bos} and \texttt{exponent}. &
\texttt{data/scalar\_zero\_modes.csv},
\texttt{data/static\_source.csv} \\
R-E31 & In the draft of the manuscript on the 4+4 spacetime the matrix
of the conserved charge was written \(W\), while the same paper writes
the kernel of the gap equation \(W(k,k')\); the main registry calls the
matrix \(W_{1644}\). The matrix \(C^-\) was used there one section
before it was defined. The draft of the manuscript on the 7+1 spacetime
distinguished the sectional curvature \(K_{AB}\) from ``the matrix
\(K\)'', which only the other paper has. And the tables of symbols of
both drafts omitted symbols that the text uses (among them \(\chi\),
\(\beta\), \(\alpha^j\), \(K_{AB}\), \(w\), \(Z_\pm\), \(\rho\),
\(\omega_b\), \(\omega_c\), \(v_F\), \(\phi\), \(W(k,k')\)). & A script
listed every letter that occurs in a formula of a manuscript and not in
its table of symbols; each was then read in its place. & \(W_{1644}\)
everywhere; \(C^-\) defined at its first use; the remark restricted to
the paper that needs it; the tables completed; the sectional curvature
and the index letters entered in the registry. & manuscripts, Sec. III
and Appendix A; RP1, Sec. 3 \\
R-E32 & The manuscripts label the four flat coordinates other than the
time by the index \(j\) and keep \(a\) for the four extra axes; DMD and
the statements of the checks label the flat coordinates by \(a\), as the
corresponding documents of the main project do. The registry of the
revision did not record the index letters at all, and the formula of the
program for the Einstein and Lanczos tensors (curvatures on pairs of
directions) was named in DMD but not written out. & Comparison of the
registry with the manuscripts when the sections of the manuscripts were
entered in it. & The registry lists the index letters with their ranges,
RP1 explains the rule, and DMD writes out the formula with the sectional
curvatures \(K_{AB}\). & RP1, Sec. 3; DMD, Sec. 1.4a \\
R-E40 & The new data files of the spectra used the column name
\texttt{tau} both for the optical depth and for the conformal time; the
bare name is retired in the revision (R-E16), and the main registry
gives \(\tau\) to the optical depth. The registry entry of the metric
perturbations listed the program name \texttt{hprime}, which is a symbol
of one of the author's notebooks. The letters \(h\), \(\eta\),
\(\alpha\), \(\theta\), \(\sigma\) of the standard reference for the
perturbation equations all have other meanings in the project. & The
checks RS3 and RS5 of the symbols failed; the letters were compared with
the registries before RS5 was written. & Columns \texttt{tau\_reio} and
\texttt{tau\_cf}; the program name removed from the registry; the
document symbols \(h_{\rm sy}\), \(\eta_{\rm sy}\), \(\alpha_{\rm sy}\),
\(\tau_{\rm cf}\) and script letters for the sums over matter (registry
group D9). & RP1; RS5, Sec. 0.2 \\
\end{longtable}

\hypertarget{errors-of-procedure}{%
\section{4. Errors of procedure}\label{errors-of-procedure}}

\begin{longtable}[]{@{}
  >{\raggedright\arraybackslash}p{(\columnwidth - 8\tabcolsep) * \real{0.0435}}
  >{\raggedright\arraybackslash}p{(\columnwidth - 8\tabcolsep) * \real{0.3043}}
  >{\raggedright\arraybackslash}p{(\columnwidth - 8\tabcolsep) * \real{0.2319}}
  >{\raggedright\arraybackslash}p{(\columnwidth - 8\tabcolsep) * \real{0.2464}}
  >{\raggedright\arraybackslash}p{(\columnwidth - 8\tabcolsep) * \real{0.1739}}@{}}
\toprule\noalign{}
\begin{minipage}[b]{\linewidth}\raggedright
No.
\end{minipage} & \begin{minipage}[b]{\linewidth}\raggedright
What was wrong
\end{minipage} & \begin{minipage}[b]{\linewidth}\raggedright
How it was found
\end{minipage} & \begin{minipage}[b]{\linewidth}\raggedright
Correction
\end{minipage} & \begin{minipage}[b]{\linewidth}\raggedright
Where to check
\end{minipage} \\
\midrule\noalign{}
\endhead
\bottomrule\noalign{}
\endlastfoot
R-E15 & A figure asked for matter of the Fermi momentum \(2\,m\), for
which no history had been computed. & The figure script failed. & The
figure shows the three densities that exist. &
\texttt{figures/fig\_matter\_history} \\
R-E24 & A commit (\texttt{d5ddbc7}) was made without running the unit
tests. One test asserts the number of exact geometry checks, which had
grown from 14 to 16, and failed. The checks themselves all passed. & The
tests were run for the provenance RP8. & The test was corrected in the
next commit; the tests are run before every commit. &
\texttt{cargo\ test} \\
R-E27 & The command \texttt{matter} contained three numbers copied by
hand from the output of the command \texttt{cosmology}: the bounds
0.0644 and 0.234 (rounded), and \(H_0t_0=0.9507\) for the age of the
model, whereas the model of the revision has \(H_0t_0=0.9516\). The last
digits of the table of drifts (RS3, Sec. 6.3) and one digit of a growth
factor (1.0590 instead of 1.0589) were affected; no conclusion was. &
Reading the program while tracing a number for the check of the numbers.
& Both commands now take the bounds and the age from one routine
(\texttt{rate\_bounds}); the tables were regenerated and the two
sentences that quoted the old digits corrected. &
\texttt{cosmo\_defl.rs}; checks \texttt{def-D4}, \texttt{def-D5};
\texttt{data/dm\_drift.csv} \\
R-E29 & The table of source files of RP8 gave the numbers of lines,
typed by hand; after later changes of the program two of them were out
of date. & The same check, which found a number that no program states.
& The table, the number of files, of lines, of unit tests and of figures
are now written by the script that writes the other tables. & RP8, Sec.
3 \\
R-E41 & Four checks of the new commands had criteria that were written
before the result was known and were not consequences of anything: that
the Hubble constant falls by less than 0.2 km/s/Mpc at the bound of
lunar laser ranging (it falls by 0.211); that the fit at the rate of
nucleosynthesis has \(\Delta\chi^2>25\) for Planck-like data (it has
11.6); that the fits of the dark-radiation families on the deflating
background differ from those of the main project by less than 1 in
\(\Delta\chi^2\) (they are systematically lower, by more); and the
statement of \texttt{def-D4} said that the sign of a change is not a
prediction where both closures give the same sign. & The checks failed,
or the statement contradicted its own numbers. & The criteria test
consistency (the fit against the Fisher matrix, one closure against the
other, the copy against the main project); the statements report the
numbers. The third case led to the joint analysis of the rate and the
extra radiation (check \texttt{def-N4}). & checks \texttt{def-C3},
\texttt{def-F4}, \texttt{def-N2}, \texttt{def-N4}, \texttt{def-D4} \\
R-E42 & The test option \texttt{quick} of the command \texttt{families}
wrote its incomplete tables over the data files of the full run. &
Noticed when the instructions of RP8 were written. & The test writes to
a temporary folder. & \texttt{cmd\_spectra.rs}; RP8, Sec. 2 \\
\end{longtable}

\hypertarget{what-the-list-shows}{%
\section{5. What the list shows}\label{what-the-list-shows}}

Forty-three errors in one day of work is not a small number, and the
reader should draw the right conclusion from it: the statements of the
revision are to be trusted to the extent that a machine check stands
behind them, and no further. Of the errors of physics, most were found
because a check with a definite expectation failed; two (R-E7, R-E23)
were found only by re-reading, and concerned the \emph{interpretation}
of a correct calculation. Two more (R-E33, R-E34) had passed every
check, because the checks tested the program against the equation that
had been written down and not the equation against the theory; they were
found only when the calculation was done again from the
eight-dimensional equations (R-E33) and with a more complete program
(R-E34). Statements of interpretation in the documents of the revision
--- in particular the verdicts of RS2 and RS3 --- deserve the reader's
independent scrutiny more than the numbers do.

The entries R-E25 to R-E29 were found while the two manuscripts were
written, when every number quoted in the running text was compared with
what the programs state. None of them changes a conclusion, but each was
a number that was typed or copied by hand where a program should have
supplied it. The comparison is now a script
(\texttt{code/scripts/check\_numbers.py}; RP8, Sec. 4), which is run
before a document is released. The last three entries (R-E30 to R-E32)
were found when the finished manuscripts were read once more, page by
page, and when their symbols were compared with their tables of symbols
and with the registry. They concern a unit and the notation, not a
result.

The entries R-E33 to R-E43 belong to the Einstein--Boltzmann
calculations of the supplement RS5. Three of them (R-E33, R-E34, R-E35)
corrected statements that earlier versions of the documents of the
revision had made; the corrections are marked in those documents where
they occur.

\end{document}
